\documentclass[11pt, a4paper]{article}
\usepackage[a4paper, margin=2.5cm]{geometry}
\usepackage[english, bidi=basic, provide=*]{babel}
\usepackage{amsmath}
\usepackage{booktabs}
\usepackage{enumitem}
\usepackage{xcolor}

\definecolor{primarycolor}{RGB}{30, 58, 138}
\definecolor{accentcolor}{RGB}{15, 118, 110}

\title{CS-433 Machine learning: project 1}
\author{Erdogan Ergit, Hugo Dutheil, Mariya Rakytyanska}
\date{}


\begin{document}
\maketitle

\section{Introduction}
For this project, we decided to implement a PCA since this method is often well adapted for classification problems such as what was given here. PCA consists of building a set principle components for each data entry in the training set, find the components of a the test set, and find the best matching components from the training set to match the output.

\section{Pipeline Architecture and Stages}

The pipeline consists of four key stages, moving from raw data preparation to final predictions:

\begin{enumerate}[leftmargin=*]
    \item \textbf{Data Cleaning and Filtering}
    \begin{itemize}
        \item \emph{Missing Value Thresholding:} The script checks both columns (features) and rows (samples/people) against a $75\%$ threshold. Any column or row containing too many missing ($\text{NaN}$) values is discarded to maintain data integrity.
        \item \emph{Imputation:} Remaining missing values in the feature matrix are replaced with zeros using `np.nan\_to\_num` to prevent numerical errors during matrix multiplication and singular value decomposition.
    \end{itemize}
    
    \item \textbf{Feature Scaling and Correlation Selection}
    \begin{itemize}
        \item \emph{Standardization:} Features are centered and scaled to have a mean of $0$ and a standard deviation of $1$, ensuring all variables contribute equally regardless of their original units.
        \item \emph{Correlation Filter:} Absolute correlations between each feature and the target variable $y$ are computed. Only the top $30\%$ of features ($70^{\text{th}}$ percentile and above) are retained to filter out noise.
    \end{itemize}
    
    \item \textbf{Dimensionality Reduction via PCA}
    \begin{itemize}
        \item \emph{Singular Value Decomposition (SVD):} PCA is implemented using SVD to transform the high-dimensional feature space into an orthogonal coordinate system (principal components).
        \item \emph{Component Selection:} The pipeline selects the top $k = 15$ principal components (`scores\_train\_k` and `scores\_test\_k`) to capture the primary variance of the dataset while ignoring minor noise.
        \item \emph{Test Set Projection:} The test dataset is scaled using the training set's mean and standard deviation, then projected into the exact same PCA coordinate space using the training rotation matrix ($\mathbf{V}^T$).
    \end{itemize}
    
    \item \textbf{Batch Processing and Nearest-Neighbor Prediction}
    \begin{itemize}
        \item \emph{Memory-Efficient Batching:} To avoid memory overflow when computing pairwise distances between large training and test sets, the test set is split into batches of size $1000$.
        \item \emph{Vectorized Euclidean Distance:} Distances are calculated efficiently using the expansion formula:
        $$\|\mathbf{v} - \mathbf{w}\|^2 = \|\mathbf{v}\|^2 + \|\mathbf{w}\|^2 - 2\langle\mathbf{v}, \mathbf{w}\rangle$$
        \item \emph{Prediction Mapping:} For each test instance, the index of the closest training sample is found using `np.argmin`, and its corresponding target value ($y$) is assigned as the prediction.
    \end{itemize}
\end{enumerate}

\section{Summary of Parameters}
\begin{table}[h]
\centering
\begin{tabular}{ll}
\toprule
\textbf{Parameter} & \textbf{Value / Choice} \\
\midrule
Missing Data Threshold & $75\%$ ($0.75$) \\
Feature Selection Percentile & Top $30\%$ ($70^{\text{th}}$ percentile cutoff) \\
Dimensionality ($k$) & $15$ principal components \\
Batch Size & $1000$ samples per iteration \\
Prediction Method & $1$-Nearest Neighbor in PCA space \\
\bottomrule
\end{tabular}
\label{tab:params}
\end{table}

\end{document}
